% T30 — «Большая таблица задач»: все задачи в одной tabular
% Колонки: №, Условие, Ответ (AnswerField). Math, класс 9, 10 задач.
\documentclass[a4paper,10pt]{article}

\usepackage{../../../../Lessons/_templates/preamble_worksheet}
\usepackage{../_shared/worksheet_check}

\usepackage{array}
\usepackage{colortbl}

\geometry{a4paper, top=12mm, bottom=12mm, left=12mm, right=12mm}

% --- Палитра: серо-стальной ---
\definecolor{wcprimary}{HTML}{2B3A55}
\definecolor{wcprimaryLight}{HTML}{C8D2E2}
\definecolor{wcprimaryBG}{HTML}{F4F6FA}

\renewcommand{\familydefault}{\sfdefault}

% --- Заголовок: тонкая полоса с номером таблицы ---
\renewcommand{\WorksheetTitle}[1]{%
  \par\noindent\begin{tikzpicture}
    \fill[wcprimary] (0,0) rectangle (\linewidth, -12mm);
    \node[anchor=west, text=white, font=\bfseries\Large\sffamily]
      at (3mm, -4mm) {ТАБЛИЦА ЗАДАЧ};
    \node[anchor=west, text=wcprimaryLight, font=\bfseries\sffamily]
      at (3mm, -9mm) {#1};
    \node[anchor=east, text=white, font=\ttfamily\small]
      at ([xshift=-3mm]\linewidth, -6mm) {N = 10};
  \end{tikzpicture}\par\vspace{4mm}}

\SetAnswerFieldSize{34mm}{8mm}

% Сокращение для строки таблицы
\newcommand{\TaskRow}[3]{% #1 номер #2 условие #3 (без аргументов, поле ответа само)
  {\bfseries\color{wcprimary}#1} & #2 & \AnswerField{#1} \\[1mm]
  \arrayrulecolor{wcprimaryLight}\hline\arrayrulecolor{black}}

\begin{document}

\WorksheetTitle{Алгебра. Подготовка к ОГЭ}
{\small\itshape\color{wcprimary!85}Класс 9 \quad·\quad T30 \quad·\quad ID: T30-table-2026-05-27}\par\vspace{2mm}
\FirstPageHeader

\noindent
\renewcommand{\arraystretch}{1.4}
\arrayrulecolor{wcprimaryLight}
\begin{tabular}{|>{\centering\arraybackslash}p{8mm}|p{105mm}|>{\centering\arraybackslash}p{40mm}|}
\hline
\rowcolor{wcprimary!85}
{\color{white}\bfseries\sffamily \#} &
{\color{white}\bfseries\sffamily Условие задачи} &
{\color{white}\bfseries\sffamily Ответ} \\\hline
\TaskRow{1}{Найди корень уравнения $5x - 14 = 2x + 1$.}{}
\TaskRow{2}{Реши неравенство $2x + 7 \le 15$. Запиши наибольшее целое решение.}{}
\TaskRow{3}{Упрости выражение $\dfrac{a^2-9}{a-3}$ при $a=11$.}{}
\TaskRow{4}{Найди значение выражения $\sqrt{169} + \sqrt{36}$.}{}
\TaskRow{5}{Реши уравнение $x^2 - 7x + 12 = 0$. Запиши меньший корень.}{}
\TaskRow{6}{Найди $25\%$ от $480$.}{}
\TaskRow{7}{Решение прогрессии: $a_1=3$, $d=4$, найди $a_{10}$.}{}
\TaskRow{8}{Вычисли $\dfrac{3}{4} : \dfrac{1}{8}$.}{}
\TaskRow{9}{Найди значение $2^7 - 2^5$.}{}
\TaskRow{10}{Реши уравнение $\dfrac{x}{3} + \dfrac{x}{6} = 9$.}{}
\end{tabular}

\WorksheetQR{T30-table/L1/v1}

\end{document}
